\chapter{Visión por Computadora Aplicada al Reconocimiento de Gestos}
\label{cp:vision-computadora}

\insertminitoc
\parindent0pt

Conocidas las características lingüísticas del \acrshort{lesho}, el siguiente paso es entender cómo
un sistema computacional puede interpretar los gestos de las manos a partir de imágenes.
Este capítulo aborda las bases técnicas de la visión por computadora aplicada al
reconocimiento de gestos: desde los fundamentos del procesamiento de imágenes hasta la
extracción y normalización de puntos de referencia anatómicos de la mano y del cuerpo, que
constituyen la entrada de los modelos de clasificación utilizados en este proyecto.

\section{Fundamentos de la Visión por Computadora}

La visión por computadora es una disciplina de la inteligencia artificial que tiene como
objetivo dotar a las máquinas de la capacidad de interpretar y comprender el contenido de
imágenes y videos digitales. Según Szeliski, esta disciplina abarca un conjunto amplio de
técnicas destinadas a analizar e interpretar imágenes, con aplicaciones que van desde la
búsqueda de imágenes y la navegación autónoma hasta tareas cotidianas de procesamiento
de fotografías y videos personales \cite{szeliski_cv_2022}.

El procesamiento de imágenes digitales constituye la base sobre la cual operan las técnicas
de visión por computadora. Una imagen digital se representa como una matriz de valores
numéricos, donde cada elemento, denominado píxel, contiene información sobre la
intensidad o el color en una posición específica. Gonzalez y Woods describen que el
procesamiento de imágenes abarca operaciones que van desde la mejora y restauración de
imágenes hasta la extracción de características y el reconocimiento de objetos, organizadas
en niveles de procesamiento de complejidad creciente \cite{gonzalez_dip_2018}.

Entre las tareas fundamentales de la visión por computadora se encuentran la detección de
objetos, que consiste en localizar la presencia y posición de entidades de interés dentro de
una imagen, y la segmentación, que separa las regiones relevantes del fondo. Estas tareas
han experimentado avances significativos con la incorporación de técnicas de aprendizaje
profundo, que han superado a los métodos tradicionales basados en características diseñadas
manualmente en numerosos dominios de aplicación \cite{szeliski_cv_2022}.

\section{Detección de Manos en Imágenes}

La detección de manos en imágenes constituye un caso particular del problema de detección
de objetos, con desafíos específicos derivados de la naturaleza articulada y deformable de
la mano. A diferencia de objetos rígidos, las manos presentan una alta variabilidad de
configuraciones, oclusiones frecuentes entre los dedos, oclusiones mutuas cuando se
emplean ambas manos, y ausencia de patrones de alto contraste que faciliten su
localización \cite{bazarevsky_blazepalm_2019}.

Los enfoques tradicionales de detección de manos se basaban en características de bajo nivel
como el color de la piel, los contornos y la forma. Estos métodos presentaban limitaciones
de robustez ante variaciones de iluminación, fondos complejos y diferencias en el tono de
piel de los usuarios. El advenimiento del aprendizaje profundo permitió superar muchas de
estas limitaciones mediante modelos capaces de aprender representaciones jerárquicas
directamente a partir de los datos \cite{szeliski_cv_2022}.

Un avance relevante en esta dirección fue la presentación por parte de Bazarevsky y Zhang,
investigadores de Google, de un detector de palma denominado BlazePalm. En lugar de
detectar la mano completa con sus dedos articulados, este modelo detecta la palma, que al
ser un objeto más rígido y de proporciones más estables, resulta significativamente más
sencillo de localizar mediante cuadros delimitadores. Este enfoque alcanza una precisión
promedio cercana al 95.7\,\% en la detección de palmas y opera en tiempo real en
dispositivos móviles \cite{bazarevsky_blazepalm_2019}.

\section{MediaPipe Hands y la Extracción de Landmarks}

Entre las soluciones disponibles para la detección y seguimiento de manos en tiempo real,
MediaPipe Hands es la herramienta seleccionada para este proyecto por su precisión,
su eficiencia en dispositivos móviles y su disponibilidad como software de código abierto.

MediaPipe es un framework de código abierto desarrollado por Google para la construcción
de canales de procesamiento de datos perceptuales de distintas modalidades, como video y
audio. El framework organiza el procesamiento como un grafo dirigido de componentes
modulares, lo que permite construir soluciones de aprendizaje automático multiplataforma
y reutilizables \cite{lugaresi_mediapipe_2019}.

La solución MediaPipe Hands implementa un canal de procesamiento compuesto por tres
modelos que operan de forma coordinada. El primero es el detector de palma BlazePalm, que
opera sobre la imagen completa y devuelve un cuadro delimitador orientado de la mano. El
segundo es un modelo de landmarks que opera sobre la región recortada definida por el
detector y devuelve puntos clave tridimensionales de alta fidelidad. El tercero es un
clasificador que asigna la configuración de puntos clave a un conjunto discreto de
gestos \cite{bazarevsky_blazepalm_2019}.

El modelo de landmarks identifica 21 puntos anatómicos por mano, cada uno con
coordenadas tridimensionales que representan su posición en el plano de la imagen y su
profundidad relativa. Para el entrenamiento de este modelo, los investigadores anotaron
manualmente alrededor de 30\,000 imágenes reales con las 21 coordenadas, y
complementaron el conjunto con datos sintéticos para mejorar la robustez del modelo ante
variaciones \cite{bazarevsky_blazepalm_2019}. Una característica relevante del canal de procesamiento es que el
detector de palma no se ejecuta en cada fotograma. La ubicación de la mano en fotogramas
sucesivos se infiere a partir de los puntos clave del fotograma anterior, y el detector solo
se reactiva cuando la confianza de la inferencia cae por debajo de un umbral determinado,
lo que optimiza el rendimiento en tiempo real \cite{bazarevsky_blazepalm_2019}.

La representación de los gestos mediante landmarks ofrece ventajas significativas frente al
procesamiento directo de píxeles. Al reducir cada mano a 21 puntos con tres coordenadas,
la dimensionalidad de la entrada disminuye drásticamente, lo que permite el uso de modelos
de clasificación más ligeros y rápidos. Además, esta representación es intrínsecamente más
robusta ante variaciones de iluminación, color de fondo y tono de piel, dado que captura la
geometría de la mano en lugar de su apariencia visual \cite{zhang_mediapipe_2020}.

\section{Normalización de Coordenadas para Invarianza Espacial}

Las coordenadas de los landmarks entregadas por MediaPipe están normalizadas respecto a
las dimensiones de la imagen, con valores en el rango de cero a uno. Sin embargo, estas
coordenadas absolutas dependen de la posición de la mano dentro del encuadre, lo que
significa que una misma seña ejecutada en distintas regiones de la imagen produce vectores
de coordenadas diferentes. Para que un modelo de clasificación reconozca una seña con
independencia de su ubicación en el encuadre, es necesario aplicar un proceso de
normalización adicional \cite{szeliski_cv_2022}.

La técnica de normalización más utilizada en el reconocimiento de gestos consiste en
expresar las coordenadas de todos los landmarks de forma relativa a un punto de referencia
fijo de la propia mano, habitualmente la muñeca. Al trasladar el origen del sistema de
coordenadas a este punto, la posición absoluta de la mano en el encuadre deja de influir en
la representación, y solo se conserva la configuración geométrica interna de la mano. Este
procedimiento confiere al sistema invarianza ante la traslación \cite{zhang_mediapipe_2020}.

De forma complementaria, suele aplicarse una normalización de escala que ajusta el tamaño
del conjunto de landmarks a un rango estándar, de modo que la distancia entre la mano y la
cámara no afecte la representación. Con ambos procedimientos, traslación y escala, la
representación de una seña se vuelve consistente independientemente de dónde y a qué
distancia se ejecute frente a la cámara, lo que mejora la capacidad de generalización del
modelo de clasificación \cite{szeliski_cv_2022}.

\section{MediaPipe Pose y la Ubicación de las Manos en el Cuerpo}

El reconocimiento de la configuración de la mano no basta para distinguir todas las señas.
Como se expuso en el capítulo anterior, la ubicación o lugar de articulación es uno de los
parámetros formacionales de las señas, por lo que una misma configuración de la mano
ejecutada en zonas distintas del cuerpo puede corresponder a señas diferentes. La
normalización respecto a la muñeca, descrita en la sección anterior, conserva la forma de la
mano pero elimina la información de dónde se realiza la seña respecto al cuerpo.

Para recuperar esa información, este proyecto incorpora MediaPipe Pose, una solución del
mismo framework que estima la postura del cuerpo a partir de la imagen. MediaPipe Pose
identifica 33 puntos de referencia corporales, entre ellos los hombros y varias anclas del
rostro como la nariz, los ojos, las orejas y la boca. A partir de estos puntos se construye
un marco de referencia del cuerpo, en el que el centro de los hombros actúa como origen y el
ancho de los hombros como unidad de escala, de modo que la posición de cada mano se expresa
en relación con el cuerpo y no con su lugar en el encuadre.

Este marco confiere a la ubicación una invarianza equivalente a la que la normalización de
la mano ofrece para la forma. Al dividir la posición de la mano entre el ancho de los
hombros, se cancela el efecto de la distancia de la persona a la cámara, ya que tanto la mano
como los hombros se ven más grandes o más pequeños en la misma proporción. La combinación de
la configuración de las manos, obtenida con MediaPipe Hands, y de su ubicación en el cuerpo,
obtenida con MediaPipe Pose, conforma la entrada del modelo de reconocimiento de señas
dinámicas que se describe en los capítulos de metodología e implementación.
